\documentclass[11pt]{article}
\usepackage[utf8]{inputenc}
\usepackage{amsmath,amssymb}
\usepackage{hyperref}
\title{Efficient Electrocatalysis Hydrogen Evolution: A Resource-Aware Approach}
\author{Anonymous}
\date{}
\begin{document}
\maketitle

\begin{abstract}
This paper studies Electrocatalysis Hydrogen Evolution. Grounded in a real, citable literature \cite{ref1,ref2,ref3,ref4,ref5,ref6,ref7,ref8},
we propose a focused method and report \textbf{illustrative} experiments.
All references below are real and verifiable (OpenAlex/arXiv); the reported
numbers are simulated to illustrate intended behaviour.
\end{abstract}

\section{Introduction}
Electrocatalysis Hydrogen Evolution is an active research area. This work targets a concrete gap identified from
the recent literature and contributes a method together with an evaluation protocol.

\section{Related Work (real literature)}
The following works, retrieved from public bibliographic APIs, frame our study:
\begin{itemize}
\item Man et al. (2011) — "Universality in Oxygen Evolution Electrocatalysis on Oxide Surfaces", ChemCatChem (cited 4666).
\item Kibsgaard et al. (2012) — "Engineering the surface structure of MoS2 to preferentially expose active edge sites for electrocatalysis", Nature Materials (cited 3214).
\item Shinagawa et al. (2015) — "Insight on Tafel slopes from a microkinetic analysis of aqueous electrocatalysis for energy conversion", Scientific Reports (cited 3205).
\item Zhu et al. (2019) — "Recent Advances in Electrocatalytic Hydrogen Evolution Using Nanoparticles", Chemical Reviews (cited 3078).
\item Voiry et al. (2013) — "Enhanced catalytic activity in strained chemically exfoliated WS2 nanosheets for hydrogen evolution", Nature Materials (cited 2618).
\item Marković (2002) — "Surface science studies of model fuel cell electrocatalysts", Surface Science Reports (cited 2141).
\end{itemize}

\section{Method}
We reduce the dominant computational cost via adaptive resource allocation, activating only the components needed per input. We instantiate this idea for Electrocatalysis Hydrogen Evolution and analyse its cost/quality trade-off.

\section{Experiments (Illustrative)}
\begin{center}
\begin{tabular}{lcc}
System & Primary Metric & Efficiency \\ \hline
Strong baseline & 0.812 & 1.00$\times$ \\
Ablation & 0.828 & 0.92$\times$ \\
\textbf{Ours} & \textbf{0.851} & \textbf{0.71$\times$}
\end{tabular}
\end{center}
\emph{Metrics simulated to illustrate the intended effect; not measured.}

\section{Conclusion}
We connected a real literature to a concrete method for Electrocatalysis Hydrogen Evolution. Future work will
validate the illustrative results with real experiments.

\begin{thebibliography}{9}
\bibitem{ref1} Isabela C. Man, Hai‐Yan Su, Federico Calle‐Vallejo et al.. Universality in Oxygen Evolution Electrocatalysis on Oxide Surfaces. \emph{ChemCatChem}, 2011. DOI:10.1002/cctc.201000397
\bibitem{ref2} Jakob Kibsgaard, Zhebo Chen, Benjamin N. Reinecke et al.. Engineering the surface structure of MoS2 to preferentially expose active edge sites for electrocatalysis. \emph{Nature Materials}, 2012. DOI:10.1038/nmat3439
\bibitem{ref3} Tatsuya Shinagawa, Angel T. Garcia‐Esparza, Kazuhiro Takanabe. Insight on Tafel slopes from a microkinetic analysis of aqueous electrocatalysis for energy conversion. \emph{Scientific Reports}, 2015. DOI:10.1038/srep13801
\bibitem{ref4} Jing Zhu, Liangsheng Hu, Pengxiang Zhao et al.. Recent Advances in Electrocatalytic Hydrogen Evolution Using Nanoparticles. \emph{Chemical Reviews}, 2019. DOI:10.1021/acs.chemrev.9b00248
\bibitem{ref5} Damien Voiry, Hisato Yamaguchi, Junwen Li et al.. Enhanced catalytic activity in strained chemically exfoliated WS2 nanosheets for hydrogen evolution. \emph{Nature Materials}, 2013. DOI:10.1038/nmat3700
\bibitem{ref6} Nina Marković. Surface science studies of model fuel cell electrocatalysts. \emph{Surface Science Reports}, 2002. DOI:10.1016/s0167-5729(01)00022-x
\bibitem{ref7} Yao Zheng, Yan Jiao, Mietek Jaroniec et al.. Advancing the Electrochemistry of the Hydrogen‐Evolution Reaction through Combining Experiment and Theory. \emph{Angewandte Chemie International Edition}, 2014. DOI:10.1002/anie.201407031
\bibitem{ref8} Ming Gong, Wu Zhou, Mon‐Che Tsai et al.. Nanoscale nickel oxide/nickel heterostructures for active hydrogen evolution electrocatalysis. \emph{Nature Communications}, 2014. DOI:10.1038/ncomms5695
\end{thebibliography}
\end{document}
